\documentclass[11pt]{article}
\usepackage{amsmath,amssymb,amsthm,hyperref}
\newtheorem{theorem}{Theorem}
\newtheorem{lemma}{Lemma}
\newtheorem{corollary}{Corollary}
\title{Prime-adic obstructions to polynomial numerator patterns in Erd\H{o}s 243}
\author{Research notes; no novelty claim}
\date{6 November 2026}
\begin{document}
\maketitle

\noindent\textbf{Status.} The full
\href{https://conjectures.io/problems/erdos243-erdos-243}{Erd\H{o}s 243 problem}
remains unsolved here. These results exclude certain exact numerator patterns,
including patterns with jumps much larger than the previously imposed
$O(\log\log C_n)$ bound. No argument makes an arbitrary counterexample's
numerator polynomial. A surviving finite modular state is not an integer
counterexample. The finite-state exclusions below are computational results
with reproducible exact verification, not Lean proofs.

Let positive integers $C_n,D_n,a_n$, with $a_n\ge2$, satisfy
\begin{equation}\label{state}
 C_{n+1}=a_nC_n-D_n,\qquad D_{n+1}=a_nD_n.
\end{equation}
In the rational-tail setting, $D_n=s\prod_{i<n}a_i$ and
$C_n=D_n\sum_{i\ge n}1/a_i$. The original asymptotic hypothesis implies
$C_{n+1}/C_n\to1$ and $\log C_n=o(n)$, but the finite modular obstructions
will not need these growth estimates.

\begin{lemma}[Persistent divisibility]\label{persistent}
If $p\mid D_N$, then $p\mid D_n$ for every $n\ge N$. If in addition
$p\mid C_m$ at some $m\ge N$, then $p\mid C_n$ for every $n\ge m$.
Consequently, if the residues of $C_n$ modulo $p$ are eventually periodic
and take both zero and a nonzero value arbitrarily late, then $p$ divides
none of the $D_n$ or $a_n$.
\end{lemma}
\begin{proof}
Both persistence assertions follow immediately from \eqref{state}.
A periodic residue pattern containing both kinds of values cannot become
permanently zero. If $p$ ever divides $a_n$, it divides $D_{n+1}$, so the
same argument applies.
\end{proof}

\begin{theorem}[Polynomial support obstruction]\label{poly}
Suppose eventually $C_n=P(n)$ for a nonconstant $P\in\mathbb Q[x]$.
Outside the finite set of primes needed to clear denominators, every prime
which ever divides $D_n$ must have the following property: the reduction
of $P$ modulo that prime is either root-free or the zero function.
If $P$ has a root modulo every sufficiently large prime, no such state
sequence exists. In particular, $P$ cannot have a rational root.
\end{theorem}
\begin{proof}
Choose $h$ with $F=hP\in\mathbb Z[x]$. Outside primes dividing $h$, the
residues of $P(n)$ are periodic. If the reduction has a root and a nonzero
value, Lemma~\ref{persistent} prohibits that prime from dividing $D_n$.
For all sufficiently large primes, a nonzero reduction of $F$ of fixed
degree cannot be the zero function. Only finitely many primes divide all
coefficients of $F$.

Under the root hypothesis, therefore, all $D_n$ and $a_n$ have prime support
in a fixed finite set $S$. Let
\[
 L_n=\operatorname{lcm}(D_0,a_0,\ldots,a_{n-1}),\qquad
 \Delta_n=D_n/L_n.
\]
The quotient is an integer. Moreover
$C_n/D_n=C_0/D_0-\sum_{i<n}1/a_i$, so multiplying by $L_n$ shows
$\Delta_n\mid C_n$. Directly,
\[
 \Delta_{n+1}/\Delta_n=\gcd(a_n,L_n).
\]
After every prime of $S$ that ever occurs has occurred, every $a_n\ge2$
has a prime factor already in $L_n$. Thus $\Delta_{n+1}\ge2\Delta_n$
eventually, contradicting the polynomial upper bound on $C_n$.
If $u/v\in\mathbb Q$ is a root of $P$, reduction modulo every sufficiently
large prime sends $u/v$ to a root, giving the final assertion.
\end{proof}

The root hypothesis is not restricted to polynomials with rational roots.
For instance $(x^2-13)(x^2-17)(x^2-221)$ has a root modulo every odd prime
other than $13,17$: if both $13$ and $17$ are nonsquares, their product is
a square. Its eventual positive values are therefore also excluded.

\begin{theorem}[An elementary quadratic exclusion]\label{quadratic}
There is no state sequence satisfying \eqref{state} with
$C_n=n^2+k$ eventually, for any integer $k$.
\end{theorem}
\begin{proof}
The parity of $n^2+k$ alternates, so Lemma~\ref{persistent} forces every
$D_n$ and $a_n$ to be odd. Choose a sufficiently late $n$ for which $C_n$
is odd, so $C_{n+1}$ is even. Eliminating $D_n$ from two steps gives
\[
 C_{n+2}=(a_{n+1}+a_n)C_{n+1}-a_n^2 C_n.
\]
Modulo four the first term vanishes and $a_n^2\equiv1$, hence
$C_{n+2}\equiv-C_n\pmod4$. But $(n+2)^2+k\equiv n^2+k\pmod4$.
An odd integer cannot equal its negative modulo four, a contradiction.
\end{proof}

\paragraph{A general finite-state obstruction.}
Suppose a proposed numerator has a known period $T$ modulo $q=p^e$,
with table $c_j$, and its residues modulo $p$ recurrently take both zero
and a nonzero value. By Lemma~\ref{persistent}, $D_n$ and $a_n$ are units
modulo $q$. Form a directed graph with vertices
\[
 (j,d),\qquad j\in\mathbb Z/T\mathbb Z,\quad d\in(\mathbb Z/q\mathbb Z)^*.
\]
For each unit $a$ satisfying $ac_j-d=c_{j+1}$ modulo $q$, include the edge
\[
 (j,d)\longrightarrow(j+1,ad).
\]
An infinite integer orbit of the proposed form induces an infinite walk
in this finite graph. An acyclic graph therefore excludes that form.
A sufficient certificate of acyclicity is an integer rank strictly
decreasing along every edge. No compatibility between different moduli
is needed when one modulus already excludes the pattern.

\paragraph{Exact completed calculations.}
The program \texttt{erdos243-polynomial-obstruction.mjs} constructs these
graphs for $c_j=j^d+k\pmod q$, exhausts every $k\pmod q$, and checks a
strictly decreasing rank on every edge of every claimed acyclic graph.
It computes every solution of the edge congruence: for
$g=\gcd(c_j,q)$ there are no solutions unless $g\mid d+c_{j+1}$;
otherwise it solves modulo $q/g$ and examines all $g$ lifts, discarding
nonunits. The completed results are:
\[
\begin{array}{c|c|c|r|r}
 p&q&\text{exponent representatives}&\text{offset cases}&\text{maximum rank}\\\hline
 2&4&2&4&3\\
 3&9&3&9&6\\
 3&27&5,11,17&81&6\\
 5&25&3,7,11,15,19&125&8\\
 11&121&13&121&18
\end{array}
\]
All $340$ graphs are acyclic. There are $1,702,554$ vertices and
$1,556,322$ edges across these computations. A separate direct enumeration
of candidate multipliers and an independent DFS cycle check agree on
every offset for the small control families. A linear polynomial control
has a modular cycle, preventing an automatic negative verdict on every
input. The output is \texttt{erdos243-polynomial-results.json}.

\begin{corollary}[Excluded power-plus-constant numerators]
If $C_n=n^d+k$ eventually, where $d\ge1$ and $k\in\mathbb Z$, then the
following necessary conditions hold:
\[
 d\ge25,\qquad d\equiv1\pmod{12},\qquad d\not\equiv13\pmod{110}.
\]
In particular every degree $1\le d\le24$ is excluded.
The same exclusions hold for a fixed positive rational multiple of
$n^d+k$.
\end{corollary}
\begin{proof}
For even $d\ge2$, reduction modulo four gives the exponent-two table.
For $d\equiv3\pmod6$, reduction modulo nine gives the exponent-three
table. For $d\equiv5\pmod6$, the unit exponent modulo $27$ is one of
$5,11,17$ modulo $18$, while a multiple of three raised to these powers
vanishes modulo $27$. Thus the listed tables cover that entire class.
Similarly, every $d\equiv3\pmod4$ reduces modulo $25$ to one of the five
listed exponents modulo $20$; nonunits vanish since $d\ge3$.
Finally $d\equiv13\pmod{110}$ gives the same table as exponent thirteen
modulo $121$, using the unit-group exponent $110$ and vanishing on
multiples of eleven.

The first four exclusions leave only $d\equiv1\pmod{12}$.
Degree one is excluded by Theorem~\ref{poly}, and degree thirteen by the
last table, so the first remaining degree is twenty-five.
For $C_n=t(n^d+k)$ with fixed positive rational $t$, set
$\widetilde D_n=D_n/t=a_n(n^d+k)-((n+1)^d+k)$ on the tail.
These are positive integers and satisfy the same recurrences with
$\widetilde C_n=n^d+k$, reducing to the preceding case.
\end{proof}

\paragraph{The full gap remains.}
No step here proves that a general numerator with $\log C_n=o(n)$ must
follow an exact polynomial. Remaining degrees are not claimed possible,
and local modular cycles are not integer orbits. The original conjecture
is neither proved nor disproved by these exclusions.
\end{document}
